% !TeX root = ../main.tex
\section{Introduction}
\label{sec:intro}

Shear wall systems are the principal lateral-force-resisting components in many high-rise residential buildings, particularly in seismic regions \cite{GB500112010Code2016,qianDesignTallBuilding2018}. The arrangement of shear walls influences stiffness distribution, load paths, and construction cost, and it must simultaneously respect architectural constraints such as functional zoning, circulation, and openings. In preliminary design, these competing requirements are commonly reconciled through iterative manual layout adjustment and repeated coordination between architectural and structural teams \cite{Steiner_2016,malaga-chuquitaypeMachineLearningStructural2022}. Such a workflow is labor-intensive and tends to restrict systematic exploration of alternative layouts.

Despite extensive design experience and code provisions, routine shear wall layout development remains inefficient for three reasons. First, the mapping from architectural drawings to structurally rational wall candidates relies on implicit, experience-driven rules that are difficult to formalize and reuse. Second, resolving spatial conflicts often requires multiple cross-discipline iterations, creating a slow feedback loop during early design. Third, the quality and consistency of outcomes can vary across designers and projects, while the limited ability to explore the design space can lead to locally satisfactory but globally suboptimal layouts. These limitations motivate automated methods that can interpret architectural constraints and generate feasible wall layouts with controllable design intent.

Recent data-driven studies have treated shear wall generation as image-to-image translation, using convolutional networks, GAN-based models, and diffusion models to map rasterized floor plans to structural layouts \cite{liaoAutomatedStructuralDesign2021a,liaoIntelligentGenerativeStructural2022a,luIntelligentStructuralDesign2022a,guIntelligentDesignShear2024,pizarroUseConvolutionalNetworks2021,zhaoIntelligentDesignShear2023,zhouStructDiffusionEndendIntelligent2024}. While pixel-based representations can capture complex spatial patterns, they introduce inefficiencies that are particularly pronounced for structural elements: thin wall lines occupy a small fraction of the image grid, topological relationships are only implicitly represented, and raster outputs require additional vectorization for downstream engineering analysis \cite{wangHighResolutionImageSynthesis2018,feiIntegratedSchematicDesign2022}.

To address these issues, graph neural networks (GNNs) have been adopted to model structural topology more explicitly \cite{Bronstein_2017,Wu_2021,Hamilton_2020,zhaoIntelligentDesignShear2023b,zhaoDesignconditioninformedShearWall2023d}. A common formulation builds graphs from geometric primitives, with nodes as wall intersections and edges as wall segments, and predicts whether each segment should be a shear wall. This representation preserves connectivity but remains low-level: individual segments carry limited information about the architectural spaces they bound. As a result, room-level constraints (e.g., functional requirements, openings, and circulation) are not directly expressed, and condition-aware control is often injected in relatively rigid ways.

A room-level representation is more aligned with the physical and architectural constraints governing shear wall placement. In typical residential plans, feasible shear wall candidates are largely restricted to room boundaries, and many practical constraints are naturally described at the room level rather than at the level of isolated wall segments. Modeling rooms as graph nodes and room adjacencies as edges therefore provides an explicit interface between architectural semantics and structural layout generation. In addition, the number of rooms per plan is usually far smaller than the number of primitive wall segments, which reduces graph scale and can improve computational efficiency. Related room-graph and room-semantics studies in architectural layout generation also support the semantic effectiveness of this abstraction \cite{nauataHouseGANRelationalGenerative2020,huGraph2PlanLearningFloorplan2020,Wang_2022}.

Based on this representation, this study formulates conditional shear wall layout generation using a GNN backbone with condition modulation. Neighborhood aggregation is implemented with an attention-based message passing network (e.g., GATv2), and conditions are integrated to guide generation toward target wall-density categories. A practical challenge is that each floor plan is typically paired with only one engineered layout, which limits direct supervision for multi-condition learning. Similar data-scarce conditional settings have motivated synthetic or counterfactual supervision strategies in other domains \cite{dasConditionalSyntheticData2022,chenCounterfactualSamplesSynthesizing2023,pavlloControllingStyleSemantics2020}. To address this data reality, training is organized into a supervised stream that learns geometric fidelity from available layouts and a counterfactual stream that enforces condition compliance under alternative conditions through density-oriented constraints.

The main contributions are as follows:

\begin{itemize}
\item A room-level graph formulation for structural layout generation that encodes architectural spaces and their adjacencies, enabling direct use of room semantics and boundary constraints while reducing the scale of graph inference relative to primitive-based formulations.
\item A conditional GNN architecture for controllable shear wall layout generation under prescribed wall-density categories, including a decoding design that predicts both wall existence and geometry.
\item A dual-stream training scheme tailored to sparse paired supervision in practice (one design per plan), combining supervised learning with counterfactual, density-oriented regularization to improve conditional compliance.
\end{itemize}

The remainder of this paper is organized as follows. Section 2 reviews related work on structural layout generation, graph-based representations in architecture-engineering-construction, and conditional generation under practical data constraints. Section 3 presents the methodology, including room-level graph construction, conditional model architecture, and training strategy. Section 4 describes the experimental setup and reports quantitative and qualitative results, including ablation analysis and case studies. Section 5 discusses implications and limitations. Section 6 concludes the paper.
